\section{Flight dynamics}
\label{sec:dynamics}

Once the machine exists its mass is fixed (\cref{as:mass}) and the problem
becomes one of differential equations. This section derives them.

\subsection{Configuration and kinematics}

The configuration space is $\SEthree = \R^3\rtimes\SOthree$, with state
$(\vect{r},R)$ and velocities $(\vect{v},\vect{\omega})$.

\begin{proposition}[Attitude kinematics]
\label{prop:kinematics}
\begin{equation}
  \dt{R} = R\,\hatmap{\vect{\omega}} .
  \label{eq:Rdot}
\end{equation}
In terms of the unit quaternion $q=(q_0,\vect{q}_v)$ representing $R$,
\begin{equation}
  \dt{q} = \tfrac12\, q \otimes (0,\vect{\omega})
         = \tfrac12
           \begin{bmatrix}
             -\vect{q}_v\transp \\
             q_0 I + \hatmap{\vect{q}_v}
           \end{bmatrix}\vect{\omega}
         \eqqcolon \tfrac12\,\Omega(\vect{\omega})\,q .
  \label{eq:qdot}
\end{equation}
\end{proposition}

\begin{proof}
Let $\vect{u}$ be a vector fixed in $\mathcal{B}$, so its inertial coordinates
are $R\vect{u}$ and its inertial time derivative is
$\vect{\omega}_{\mathcal I}\times R\vect{u}$ where
$\vect{\omega}_{\mathcal I}=R\vect{\omega}$. Then
$\dt{R}\vect{u} = \hatmap{(R\vect{\omega})}R\vect{u}$ for all $\vect{u}$, so
$\dt{R} = \hatmap{(R\vect{\omega})}R = R\hatmap{\vect{\omega}}R\transp R
= R\hatmap{\vect{\omega}}$ using \cref{eq:hatid}. \Cref{eq:qdot} follows by
differentiating \cref{eq:quat2rot} and using quaternion multiplication;
see~\cite[\S1.4]{markley2014}.
\end{proof}

\begin{remark}
\Cref{eq:qdot} preserves $\|q\|$ exactly, since
$q\transp\dt q = \tfrac12 q\transp\Omega q = 0$ by antisymmetry. Numerically,
integration error accumulates and the engine renormalises after each step; see
\cref{sec:numerics}.
\end{remark}

\subsection{Newton-Euler equations}

\begin{theorem}[Equations of motion]
\label{thm:newtoneuler}
For a rigid body of mass $m$ and body-frame inertia tensor $J$ subject to a
total external force $\vect{F}$ expressed in $\mathcal I$ and a total external
moment $\vect{\tau}$ about the centre of mass expressed in $\mathcal B$,
\begin{align}
  \dt{\vect{r}} &= \vect{v}, \label{eq:rdot}\\
  m\,\dt{\vect{v}} &= \vect{F}, \label{eq:vdot}\\
  J\,\dt{\vect{\omega}} + \vect{\omega}\times J\vect{\omega}
    &= \vect{\tau} . \label{eq:wdot}
\end{align}
\end{theorem}

\begin{proof}
\Cref{eq:rdot,eq:vdot} are Newton's second law in an inertial frame. For
\cref{eq:wdot}, the angular momentum about the centre of mass expressed in
$\mathcal I$ is $\vect{h}_{\mathcal I} = RJ\vect{\omega}$ and
$\dt{\vect h}_{\mathcal I} = \vect{\tau}_{\mathcal I} = R\vect\tau$.
Differentiating using \cref{eq:Rdot},
\[
  \dt{\vect h}_{\mathcal I}
  = \dt R J\vect\omega + RJ\dt{\vect\omega}
  = R\hatmap{\vect\omega}J\vect\omega + RJ\dt{\vect\omega}
  = R\big(\vect\omega\times J\vect\omega + J\dt{\vect\omega}\big).
\]
Equating and cancelling $R$ gives \cref{eq:wdot}.
\end{proof}

The forces are
\begin{equation}
  \vect{F} = T\,R\ez - mg\,\ez + \vect{F}_{\mathrm{aero}}
             + \vect{F}_{\mathrm{dist}} ,
  \label{eq:forces}
\end{equation}
with $T$ the total rotor thrust along the body $\vect b_3$ axis.

\subsection{Rotor wrench and the allocation map}

\begin{model}[Quasi-steady rotor]
\label{mod:rotorcoeff}
Rotor $i$ turning at $\Omega_i$ produces thrust and reaction torque
\begin{equation}
  T_i = k_T\Omega_i^2, \qquad Q_i = k_Q\Omega_i^2 ,
  \label{eq:ktkq}
\end{equation}
with $k_T,k_Q$ determined by \cref{sec:momentum,sec:blade} rather than fitted:
$k_T = T/\Omega^2$ and $k_Q = \Pshaft/\Omega^3$ evaluated at the design point.
\end{model}

Let rotor $i$ sit at body position $\vect{d}_i = (x_i,y_i,0)$ and have spin
sense $s_i\in\{+1,-1\}$. Its contribution to the body moment is
$\vect{d}_i\times T_i\ez$ from thrust and $-s_iQ_i\ez$ from reaction torque, so
\begin{equation}
  \vect\tau = \sum_{i=1}^{N}
    \begin{bmatrix} y_i\,k_T \\ -x_i\,k_T \\ -s_i\,k_Q \end{bmatrix}\Omega_i^2 ,
  \qquad
  T = \sum_{i=1}^{N} k_T\Omega_i^2 .
  \label{eq:wrench}
\end{equation}

\begin{definition}[Allocation matrix]
Collecting \cref{eq:wrench} with $\vect{u} = (\Omega_1^2,\dots,\Omega_N^2)$,
\begin{equation}
  \begin{bmatrix} T \\ \vect\tau \end{bmatrix}
  = \mat{M}\,\vect{u},
  \qquad
  \mat{M} =
  \begin{bmatrix}
    k_T & \cdots & k_T\\
    y_1 k_T & \cdots & y_N k_T\\
    -x_1 k_T & \cdots & -x_N k_T\\
    -s_1 k_Q & \cdots & -s_N k_Q
  \end{bmatrix} \in \R^{4\times N} .
  \label{eq:alloc}
\end{equation}
\end{definition}

\begin{proposition}[Controllability of the wrench]
\label{prop:alloc}
The set of achievable wrenches is $\mat{M}\R_{\ge0}^N$. Full four-axis
authority requires $\rank \mat M = 4$, which for hubs equispaced on a circle of
radius $L$ with alternating spin senses holds whenever $N\ge4$ is even. For odd
$N$, $\sum_i s_i\ne0$ and the hover trim $\Omega_i\equiv\Omega_h$ produces a
nonzero yaw moment, so the vehicle cannot trim without unequal thrusts.
\end{proposition}

\begin{proof}
The rank statement is a direct computation on \cref{eq:alloc}: the first row is
constant, rows two and three span the $(y_i,-x_i)$ pattern which has rank two
for $N\ge3$ equispaced hubs, and row four is independent of the others exactly
when the $s_i$ are not all equal. For odd $N$ with alternating signs,
$\sum_i s_i = \pm1$, so at equal $\Omega_i$ the yaw component of
$\mat M\vect u$ is $\mp k_Q\Omega_h^2 \ne 0$.
\end{proof}

\begin{remark}
\Cref{prop:alloc} is why the engine restricts $N$ to even values. A tricopter
resolves the trim problem with a tilting motor, which is a different actuation
model and is outside the scope of \cref{eq:alloc}.
\end{remark}

\subsection{Motor dynamics}

Rotor speed does not follow command instantaneously. The simplest adequate
model is first order,
\begin{equation}
  \tau_m\,\dt{\Omega}_i = \Omega_i^{\mathrm{cmd}} - \Omega_i ,
  \label{eq:motorlag}
\end{equation}
with $\tau_m$ of order \SIrange{20}{60}{\milli\second}. The underlying
electromechanical model is
\begin{equation}
  L\dt{\imath}_i = V_i - R_m \imath_i - k_e\Omega_i,
  \qquad
  J_r\dt{\Omega}_i = k_\tau \imath_i - k_Q\Omega_i^2 ,
  \label{eq:bldc}
\end{equation}
whose electrical time constant $L/R_m$ is typically two orders of magnitude
faster than the mechanical one, so that singular perturbation reduces
\cref{eq:bldc} to a first-order system in $\Omega_i$; \cref{eq:motorlag} is its
linearisation about the trim point. We use \cref{eq:motorlag} and record the
simplification.

\subsection{Battery state}

\begin{equation}
  \dt{E} = -P, \qquad
  P = \frac{1}{\eta}\sum_{i=1}^{N} k_P\,\Omega_i^{3} + \Paux ,
  \qquad
  k_P = \frac{\Pshaft}{\Omega^3}\bigg|_{\text{design}} ,
  \label{eq:edot}
\end{equation}
or in normalised form $\dt{s} = -P/E_{\max}$ with $s = E/E_{\max}$ the state of
charge. The cubic dependence follows from $\Pshaft = Q\Omega$ with
$Q\propto\Omega^2$. \Cref{eq:edot} is the coupling that makes control effort
cost endurance, and it is what \cref{sec:codesign} closes.

\Cref{eq:edot} is calibrated at hover and has no inflow term. At fixed rotor
speed it returns the same power in hover and in axial translation, so it does
not contain the climb correction of \cref{eq:axial}, which the component model
of \cref{sec:momentum} does. It is a near-hover approximation, and so are the
mission energies computed with it: they are appropriate for following a person
indoors at walking speed and vertical rates of a few tenths of a metre per
second, and they are not a model of climbs, descents or aggressive manoeuvres.
Adding a climb term to the power alone, while thrust keeps its hover
calibration, would make the two inconsistent; a consistent treatment couples an
axial and oblique inflow model to both.

\subsection{The screen as an aerodynamic surface}

This is what distinguishes the vehicle from a camera drone. The display is a
large flat plate, rigidly attached through a gimbal, and therefore both
generates orientation-dependent drag and applies that drag at a point offset
from the centre of mass.

\begin{model}[Flat plate with orientation]
\label{mod:screen}
Let $\vect{n}$ be the outward normal of the screen in inertial coordinates and
$\vect{v}_{\mathrm{rel}} = \vect{v} - \vect{v}_{\mathrm{air}}$. Decompose
$\vect{v}_{\mathrm{rel}} = v_n\vect{n} + \vect{v}_t$ with
$v_n = \vect{n}\transp\vect{v}_{\mathrm{rel}}$. Then
\begin{equation}
  \vect{F}_{\mathrm{screen}}
  = -\tfrac12\rho A_s\Big[
      C_{D,n}\,|v_n|\,v_n\,\vect{n}
      + C_{D,t}\,\|\vect{v}_t\|\,\vect{v}_t \Big] ,
  \label{eq:screendrag}
\end{equation}
with $C_{D,n}\approx1.17$ for a normal flat plate and $C_{D,t}\ll C_{D,n}$
skin friction.
\end{model}

The moment about the centre of mass follows from the offset $\vect{d}_s$ of the
screen's centre of pressure in body coordinates:
\begin{equation}
  \vect\tau_{\mathrm{screen}} = R\transp\big[(R\vect{d}_s)\times
    \vect{F}_{\mathrm{screen}}\big] .
  \label{eq:screenmoment}
\end{equation}

\begin{remark}[The sign of $\vect{d}_s$ matters]
If the display is carried \emph{above} the rotor plane, $d_{s,3}>0$; if it
hangs below, $d_{s,3}<0$. In forward flight the drag acts rearward, so a screen
above the centre of mass produces a nose-up moment that opposes the forward
tilt, a weathercock-stable arrangement, while a screen below produces a
nose-down moment that reinforces it. The magnitude is identical; only the sign
changes, and it changes the closed-loop behaviour measurably.
\end{remark}

\begin{remark}[Why the screen is nearly transparent to the rotors]
The screen is vertical and the rotor inflow and wake are vertical, so the plate
is edge-on to the rotor flow and blocks a projected area of order
(thickness $\times$ width), some \SI{2e-3}{\metre\squared} against a disk area
of \SI{0.76}{\metre\squared}. Placing it above or below the rotor plane is
therefore almost free aerodynamically. What it changes is geometry relative to
the user, which \cref{sec:results} shows to be the decisive consideration.
\end{remark}

\subsection{Gimbal}

The display is carried on a gimbal whose job is to hold it upright while the
airframe tilts to accelerate. Modelling the pitch and roll axes as
independently actuated single-degree-of-freedom systems,
\begin{equation}
  J_g\ddt{\theta}_{g,j} + b_g\dt{\theta}_{g,j}
    = \tau_{g,j} - \tau_{\mathrm{wind},j},
  \qquad
  |\tau_{g,j}| \le \tau_g^{\max},
  \quad j\in\{\text{pitch},\text{roll}\} .
  \label{eq:gimbal}
\end{equation}
The gimbal motor torque reacts on the airframe and is included in $\vect\tau$.

\begin{definition}[Readability error]
Let $\vect{u}_s$ be the unit vector along the top edge of the display in
inertial coordinates. The readability error is
\begin{equation}
  \theta_{\mathrm{read}} = \arccos\big(\ez\transp\vect{u}_s\big) ,
  \label{eq:readability}
\end{equation}
the angle between the top of the image and vertical.
\end{definition}

\begin{remark}
\Cref{eq:readability} is deliberately not the angle between the screen normal
and horizontal. Gimbal roll leaves the normal unchanged while rotating the
image in its own plane, so a normal-based metric is blind to exactly the error a
reader notices most. Measuring the up-vector captures both axes.
\end{remark}

\subsection{The complete system}

Assembling, the state is
\begin{equation}
  x = \big(\vect{r},\ \vect{v},\ q,\ \vect{\omega},\
        \Omega_1,\dots,\Omega_N,\ E,\
        \vect\theta_g,\ \dt{\vect\theta}_g\big)
  \in \R^{3}\times\R^{3}\times\mathbb{S}^{3}\times\R^{3}\times\R^{N}\times\R
      \times\R^{2}\times\R^{2} ,
  \label{eq:state}
\end{equation}
of dimension $18+N$, and the dynamics are
\begin{equation}
  \dt{x} = f(x,u,w;\,p),
  \label{eq:ode}
\end{equation}
with control $u = (\Omega^{\mathrm{cmd}}_{1:N}, \vect\tau_g)$, disturbance $w$
comprising wind and human motion, and design parameters $p$.

Crucially, \cref{eq:ode} is not the whole problem. The parameters $p$ are
themselves constrained by the algebraic closure of \cref{sec:closure}, giving a
differential-algebraic system
\begin{equation}
  \dt{x} = f(x,u,w;p),
  \qquad
  0 = h(x,u,p) ,
  \label{eq:dae}
\end{equation}
whose algebraic part contains \cref{eq:budget} with the mission energy computed
from a solution of the differential part. \Cref{sec:codesign} makes this
precise.
